\documentclass[10pt,a4paper]{amsart}
\usepackage[margin=19mm]{geometry}
\usepackage{amsmath,amssymb,mathtools}
\usepackage[scr=rsfs,cal=euler]{mathalfa}
\usepackage{xfrac}
\usepackage[unicode,hidelinks,bookmarks=false]{hyperref}
\newcommand{\ie}{i.e.\ }
\newcommand{\cf}{cf.\ }
\newcommand{\resp}{respectively\ }
\newcommand{\Subsection}{\subsection}
\providecommand{\nobreakdash}{\nobreak\hbox{-}}
\setlength{\emergencystretch}{2em}
\multlinegap0pt
\usepackage{enumerate}
\usepackage{algorithm}\usepackage{algpseudocode}
\usepackage{xcolor}
\usepackage{tikz}
\newtheorem{lemma}{Lemma}
\title{Sparsity greedoids and pebble game algorithms for posets}
\author{Signe Lundqvist, Tovohery Randrianarisoa, Klara Stokes, Joannes Vermant}
\date{Source: arXiv:2306.05050v3 (CC BY 4.0)}
\begin{document}
\maketitle

\noindent\emph{Source and licence.} Sparsity greedoids and pebble game algorithms for posets. Version arXiv:2306.05050v3 (CC BY 4.0), distributed by arXiv under the Creative Commons Attribution 4.0 International licence, http://creativecommons.org/licenses/by/4.0/, as recorded in the arXiv licence metadata for this exact version and shown on the abstract page https://arxiv.org/abs/2306.05050v3. Full licence notice: \texttt{licenses/arxiv-2306.05050-CC-BY-4.0.txt}.\par\smallskip

\noindent\textit{Editorial scope.} Counted unit: the article's lemma surplus-independent (source0.tex lines 676--681). The article's complete original proof of it is delivered with it (source0.tex lines 682--726, one \texttt{proof} environment of 45 lines). No mathematics is written, repaired or re-derived: the statement and the proof are the article's own lines, in the article's own notation, and no retained line is substituted or re-flowed.  Retained uncounted as the source-local context the proof needs: lines 241--245; lines 600--611; lines 631--640.  Nothing was omitted from any retained span, so the package carries no omission record.  No retained line contains a citation, so the wrapper carries no bibliography and no bibliography entry.  No retained line contains a figure environment, an image-inclusion command or drawing code, so no image is delivered and the package declares no asset dependency.  Editorial formatting only, in addition: an AMS \texttt{amsart} preamble that restates the article's own declarations and macros used by the retained lines, namely the environments \texttt{lemma} on separate counters, together with the standard packages those lines need.  The retained environments print in this document as Lemma 1 in order of appearance, whereas the article prints the same items with the numbering of its own full document.  The complete native source of the article as distributed in the arXiv e-print for this exact version is bundled unchanged as \texttt{sources/142659/source0.tex}.
\begin{align}
    k_i + k_{i+1} &> l_{i, i+1} \label{conditionposet1 -k k l},\\
    k_{i+1} + l_{i-1, i} &> l_{i, i+1}\label{conditionposet2 - k l l},\\
     l_{1,2} \leq  \dots &\leq l_{n-1,n}. \label{conditionposet3-increasing}
\end{align}

\begin{lemma}\label{Increase_when_tight - greedoids}
Let $K\in \mathbb{Z}_{>0}^{n}, L\in \mathbb{Z}_{\geq 0}^{n-1}$ be such that inequalities \eqref{conditionposet1 -k k l}, \eqref{conditionposet2 - k l l} and \eqref{conditionposet3-increasing} are valid. Let $G=(V,E)$ be a graph such that the vertex set has a partition $V=V_1\cup \dots \cup V_n$ such that all edges $e\in E$ have endpoints in $V_{i} \cup V_{i+1}$ for some $1 \leq i < n$. 

Suppose also that $B \subseteq E$ is $(K,L)$-sparse. Let $X\subseteq V_{1}\cup \dots \cup V_{j}$ and $T \subseteq V_{1}\cup \dots \cup V_{j}$ be such that $X\cap T$ contains some vertices of $V_{j-1}$ and $V_{j}$. Suppose also that
\begin{equation*}
\sum_{i=1}^{j} k_i |V_i\cap T| - |B(T)|  = \ell_{j-1, j}. 
\end{equation*}
Then
\begin{equation*}
\sum_{i=1}^{j} k_i |V_i \cap X| - |B(X)|  \geq \sum_{i=1}^{j} k_i |V_i\cap (T\cup X)| - |B(T\cup X)|. 
\end{equation*}
\end{lemma}

\begin{lemma}\label{blocks}
Let $K\in \mathbb{Z}_{>0}^{n}, L\in \mathbb{Z}_{\geq 0}^{n-1}$ be such that inequalities \eqref{conditionposet1 -k k l}, \eqref{conditionposet2 - k l l} and \eqref{conditionposet3-increasing} are valid. Let $G=(V,E)$ be a graph such that the vertex set has a partition $V=V_1\cup \dots \cup V_n$ such that all edges $e\in E$ have endpoints in $V_{i} \cup V_{i+1}$ for some $1 \leq i < n$, and assume that $G$ is $(K,L)$-sparse. 

Let $B_1, B_2 \subseteq E$ be tight edge sets, with $V(B_i)\subseteq V_1\cup \dots \cup V_t$. If $V(B_1)\cap V(B_2) \cap V_{t-1} \neq \emptyset$ and $V(B_1)\cap V(B_2) \cap V_{t} \neq \emptyset$, then $B_1 \cap B_2$ and $B_1\cup B_2$ are tight subgraphs.
Moreover, one has
\begin{align*}
    V(B_1\cup B_2) &= V(B_1) \cup V(B_2)\\
    V(B_1\cap B_2) &= V(B_1) \cap V(B_2)
\end{align*}
\end{lemma}

\begin{lemma}
    Let $K\in \mathbb{Z}_{>0}^{n}, L\in \mathbb{Z}_{\geq 0}^{n-1}$ be such that inequalities \eqref{conditionposet1 -k k l}, \eqref{conditionposet2 - k l l} and \eqref{conditionposet3-increasing} are valid. Let $G=(V,E)$ be a graph such that the vertex set has a partition $V=V_1\cup \dots \cup V_n$ such that all edges $e\in E$ have endpoints in $V_{i} \cup V_{i+1}$ for some $1 \leq i < n$. 
    
    If $\mathcal{M}_c$ is a matroid for all $1 \leq c\leq j$, then $\theta_{\alpha_j} = \theta_{\alpha_j'}$ for every two words $\alpha_j$ and $\alpha_j'$, such that $E_{\alpha_j}$ and $E_{\alpha_j'}$ are maximal feasible subsets supported on $V_1\cup \dots \cup V_{j+1}$. 
    \label{surplus-independent}
\end{lemma}

\begin{proof}
We are given $\alpha_j$ and $\alpha_j'$. For any word $\alpha_j = e_1, \dots e_m$ such that $E_{\alpha_j}$ is supported on $V_1\cup \dots \cup V_{j+1}$ and $d<j$, define $\alpha_{d}$ to be the word $e_1 \dots e_{\tilde{m}}$ obtained by restricting to those edges supported on $V_1\cup \dots \cup V_{d+1}$. I.e. $\alpha_{d}$ is the word resulting from deleting the edges $e_{\tilde{m}+1}\dots e_{m}$ where $\tilde{m}+1$ is the smallest index such that for all $i\geq \tilde{m}+1$, one has $e_{i} \in E(V_{d+i}\cup V_{d+i+1})$ for some $i \geq 1$. In particular, $\alpha_{0}$ will always be the empty word and for subsets $X\subseteq V_1$, that $\theta_{\alpha_0}(X) = \theta_{\alpha_0'}(X) = k_1|X|$, which serves as the basis for our induction. Suppose now that the statement holds for all $d < j$. We will show it holds for $d=j$.

Let $X\subseteq V_{j+1}$. One can verify that one can express $\theta_{\alpha_{j}}$ as a function of $\theta_{\alpha_{j-1}}$. Namely, one has:
\begin{equation*}
    \theta_{\alpha_j}(X)= \min_{\substack{W\subseteq V_{j} \cup V_{j+1}\\[2pt] X\subseteq W}}  k_{j+1}|V_{j+1} \cap W| +\theta_{\alpha_{j-1}}(V_{j}\cap W)  -|E_{\alpha_j}(W)|.
\end{equation*}


By induction, we can assume that $\theta_{\alpha_{j-1}} = \theta_{\alpha_{j-1}'}$. Thus, it suffices to consider the case where $\alpha_{j}$ and $\alpha_{j}'$ only differ for edges $e\in E(V_{j}\cup V_{j+1})$, and we may assume $\alpha_{j}\neq \alpha_{j}'$ since otherwise the lemma is trivially true. In this case, since we know that $\mathcal{M}_{j}$ is a matroid, we only need to show that in the case where $E_{\alpha_{j}} = \left (E_{\alpha_{j}'} \cup \{e\}\right) \setminus \{f\}$, one has $\theta_{\alpha_{j}}(X) = \theta_{\alpha_{j}'}(X)$, since the statement then follows for all maximally chosen words $\alpha_j$ and $\alpha_j'$. Suppose that the added edge, $e$, has endpoints $v,w$. By Lemma \ref{blocks}, we know that the intersection 
\begin{equation*}
M  = \bigcap_{\substack{T\subseteq E_{\alpha_j'},\\ v,w \in V(T),\\ T \textup{ tight }}} T
\end{equation*}
is a nonempty tight subset of $E_{\alpha_j'}$.  Furthermore, if $f$ were not in $M$, then $M$ would be a tight subset of $E_{\alpha_j}$, as $f$ is the only edge of $E_{\alpha_j'}$ not in $E_{\alpha_j}$. Then, $M \cup \{e\}$ would not be $(K,L)$-sparse, which contradicts $E_{\alpha_j}$ being $(K,L)$-sparse. It follows that $f \in M$.  Now, let $X\subseteq V$, and suppose $W$ minimises the expression for $\theta_{\alpha_j'}$, so
\begin{equation*}
\theta_{\alpha_{j}'}(X) = \sum_{i=1}^{j+1} k_i|V_i\cap W| -|E_{\alpha_j'}(W)|.
\end{equation*}

 If $f \notin E_{\alpha_j'}(W)$, then we see that

 \[|E_{\alpha_j}(W)| \geq |E_{\alpha_j'}(W)|\]

so it follows immediately that 

\begin{align*}
    \theta_{\alpha_j}(X) &\leq  \sum_{i=1}^{j+1} k_i |V_i \cap W| - |E_{\alpha_{j}}(W)| \\
    &\leq \sum_{i=1}^{j+1} k_i |V_i \cap W| - |E_{\alpha_{j}'}(W)| = \theta_{\alpha_j'}(X).
\end{align*}
 
 If $f \in E_{\alpha_j'}(W)$, then we get using Lemma \ref{Increase_when_tight - greedoids}, that
\begin{align*}
\sum_{i=1}^{j+1} k_i |V_i\cap W| - |E_{\alpha_j'}(W)|  \geq \sum_{i=1}^{j+1} k_i |V_i\cap (W\cup V(M))| - |E_{\alpha_{j}'}(W \cup V(M))|,
\end{align*}
and hence we may assume that $V(M) \subseteq W$.

Since $E_{\alpha_{j}} = \left (E_{\alpha_{j}'} \cup \{e\}\right) \setminus \{f\}$, and $e$ and $f$ are both supported on $V(M)$ it follows that,
$$|E_{\alpha_{j}'}(W\cup V(M))|=|E_{\alpha_{j}}(W\cup V(M))|.$$

It now follows that

$$\theta_{\alpha_j}(X) \leq  \sum_{i=1}^{j+1} k_i |V_i \cap W| - |E_{\alpha_{j}'}(W)| = \theta_{\alpha_j'}(X).$$


We conclude that $\theta_{\alpha_j}(X)\leq \theta_{\alpha_j'}(X)$ in all cases. Switching the roles of $E_{\alpha_j}$ and $E_{\alpha_j'}$, the same argument gives that $\theta_{\alpha_j}(X) \geq \theta_{\alpha_j'}(X)$, so $\theta_{\alpha_j'}(X) = \theta_{\alpha_j}(X)$. 
\end{proof}

\end{document}
